\section{Experiments and Results}
\label{sec:experiments}

To rigorously evaluate our proposed \textbf{PFSR} method and the \textbf{OTSP} extension, we conduct extensive experiments in a high-fidelity calibrated representation sandbox. We benchmark them against multiple state-of-the-art ensembling baselines across three deployment configurations and analyze their performance, stability, and susceptibility to overfitting.

\subsection{Experimental Design and Setup}
Our experimental environment simulates task representations in a $D=192$ dimensional feature space across $K=4$ diverse classification tasks (MNIST, FashionMNIST, CIFAR-10, and SVHN), representing a standard multi-task expert registry. To simulate realistic Vision Transformer embeddings, representations are generated within task-specific orthogonal subspaces with task-specific noise scales ($\sigma = [0.05, 0.25, 0.40, 1.95]$) and subspace-isolated noise to prevent random leakage. We instantiate task-specific expert classifiers perfectly aligned with the orthogonal class prototypes to serve as stable decision boundaries, establishing a joint oracle ceiling accuracy of \textbf{74.46\%}. All results are averaged over 10 independent random seeds (seeds 42 to 51) to ensure strict statistical significance.

We evaluate and compare the following methods:
\begin{itemize}
    \item \textbf{Uniform Merging (Task Arithmetic):} Parameter-free baseline setting static uniform ensembling weights ($\alpha_k = 0.25$) for all samples.
    \item \textbf{Parametric LinearRouter:} A standard single-layer Linear Router trained on a 64-sample calibration split via AdamW optimization.
    \item \textbf{QWS-Merge (SOTA Parametric):} An over-parameterized wave-superposition routing layer trained with weight decay.
    \item \textbf{L3-Softmax (Unregularized):} A trainable Softmax router with random weight initialization.
    \item \textbf{L3-Softmax Well-Reg (Zero-Init):} A trainable Softmax router with zero-initialization and weight decay.
    \item \textbf{PFSR (Ours, Parameter-Free):} Our proposed parameter-free projection onto unorthogonalized task centroids.
    \item \textbf{OTSP (Ours, Orthogonal):} Our proposed orthogonalization extension utilizing L{\"o}wdin Symmetric Orthogonalization to build an orthonormal basis.
\end{itemize}

\begin{table*}[t]
\caption{Joint Mean Accuracy (Mean $\pm$ Standard Deviation \%) and Routing Accuracy (\%) across 10 random seeds under three deployment configurations: Homogeneous ($B=256$), Heterogeneous ($B=256$), and Heterogeneous ($B=1$). Under uncorrupted (clean centroid) representations, our training-free PFSR and OTSP achieve perfect 100\% routing accuracy, while the unconstrained LinearRouter baseline exhibits instability under $B=1$ vectorized streaming.}
\label{tab:main_results}
\vskip 0.15in
\begin{center}
\begin{small}
\begin{sc}
\setlength{\tabcolsep}{3.5pt}
\begin{tabular}{lcccc}
\toprule
Router Method & Homogeneous & Heterogeneous & Heterogeneous & Routing \\
& ($B=256$) & ($B=256$) & ($B=1$) & Accuracy \\
\midrule
\textbf{Expert Ceiling Reference} & \textbf{74.46\% $\pm$ 0.81\%} & \textbf{74.46\% $\pm$ 0.81\%} & \textbf{74.46\% $\pm$ 0.81\%} & -- \\
\midrule
Uniform Merging & 74.46\% $\pm$ 0.81\% & 74.46\% $\pm$ 0.81\% & 74.46\% $\pm$ 0.81\% & 25.00\% $\pm$ 0.00\% \\
Naive Mean Centroid (Symmetric Collapse) & 74.46\% $\pm$ 0.81\% & 74.46\% $\pm$ 0.81\% & 74.46\% $\pm$ 0.81\% & 25.18\% $\pm$ 1.10\% \\
Parametric LinearRouter (Unreg) & 72.40\% $\pm$ 2.10\% & 70.75\% $\pm$ 1.97\% & 55.57\% $\pm$ 1.68\% & 66.65\% $\pm$ 1.53\% \\
QWS-Merge (SOTA Parametric) & 74.46\% $\pm$ 0.81\% & 74.46\% $\pm$ 0.81\% & 74.46\% $\pm$ 0.81\% & 51.24\% $\pm$ 3.72\% \\
L3-Softmax (Unregularized) & 74.46\% $\pm$ 0.81\% & 74.46\% $\pm$ 0.81\% & 74.46\% $\pm$ 0.81\% & 66.47\% $\pm$ 1.62\% \\
L3-Softmax Well-Reg (Zero-Init) & 74.46\% $\pm$ 0.81\% & 74.46\% $\pm$ 0.81\% & 74.46\% $\pm$ 0.81\% & 67.22\% $\pm$ 1.61\% \\
\textbf{PFSR (Ours, Parameter-Free)} & \textbf{74.46\% $\pm$ 0.81\%} & \textbf{74.46\% $\pm$ 0.81\%} & \textbf{74.46\% $\pm$ 0.81\%} & \textbf{100.00\% $\pm$ 0.00\%} \\
\midrule
OTSP (Orthogonal Extension) & 74.46\% $\pm$ 0.81\% & 74.46\% $\pm$ 0.81\% & 74.46\% $\pm$ 0.81\% & 100.00\% $\pm$ 0.00\% \\
\bottomrule
\end{tabular}
\end{sc}
\end{small}
\end{center}
\vskip -0.1in
\end{table*}

\subsection{Quantitative Results and Discussion}
The comprehensive results of our 10-seed sweep are detailed in Table~\ref{tab:main_results}. The empirical findings yield profound qualitative and methodological insights.

\textbf{1. Perfect Routing Stability, Naive Centroid Collapse, and the Orthogonal Masking Effect.}
Under uncorrupted expert weight matrices (where non-target class coordinates are properly initialized to zero), our non-parametric dynamic routing methods (**PFSR** and **OTSP**) achieve a perfect \textbf{100.00\% $\pm$ 0.00\%} routing accuracy. This flawless performance empirically validates our mathematical formulation of task centroids via Singular Value Decomposition (SVD): extracting the first principal component of each specialist's weight matrix captures the true, uncorrupted direction of maximum task variance.

To empirically demonstrate the necessity of SVD centroid extraction over naive class prototype averaging, we evaluate the \textbf{Naive Mean Centroid} baseline, which routes samples by projecting them onto the simple average of each expert's class prototype weights. As shown in Table~\ref{tab:main_results}, because the classifier weights are symmetrically distributed around the origin (sum-to-zero) to maximize margins, the simple average of the prototype vectors collapses to near $\mathbf{0}$, forcing its routing accuracy to drop to near-random guessing (\textbf{25.18\% $\pm$ 1.10\%}). In stark contrast, our proposed SVD-based centroid extraction captures the principal direction of maximum variance, completely avoiding sum-to-zero cancellation and maintaining perfect 100.00\% routing accuracy.

However, Table~\ref{tab:main_results} also illustrates a key methodological insight regarding sandbox evaluation. Because the task representations reside in perfectly disjoint orthogonal subspaces, any positive ensembling weight assigned to the correct expert returns the ceiling accuracy. Consequently, the joint classification accuracy is completely flat and matches the Expert Ceiling Reference (\textbf{74.46\% $\pm$ 0.81\%}) across Uniform, QWS-Merge, L3-Softmax, PFSR, and OTSP. This is a direct consequence of the \textit{Orthogonal Masking Effect}, where multiplying the correct task's logits by any positive scalar is argmax-invariant because the out-of-subspace logits of all other experts are exactly $0.0$. Thus, joint classification accuracy in perfectly disjoint setups is insensitive to routing quality, making routing accuracy the only sensitive benchmark.

\textbf{2. Robust parametric baseline training and the small-sample overfitting penalty.}
To establish a mathematically rigorous, fully capacity-optimized, and fair comparison, we updated the training methodology of the parametric baselines. Rather than using batch-averaged ensembling weights, we train all parametric routers (LinearRouter, QWS-Merge, L3-Softmax, L3-Softmax Well-Reg) directly on the raw, $D$-dimensional representations of the calibration split to predict the true task labels $task_{cal}$ via a standard cross-entropy loss. This provides a direct, sample-wise supervised training objective that represents the optimal training configuration for parametric classifiers.

Even under this highly optimized setting, our sweep reveals that while parametric routers achieve a routing accuracy of **51.24\% to 67.22\%** on the test split, they are highly susceptible to small-sample inductive overfitting on the tiny calibration split (only 16 samples per task, 64 samples total). In this low-data regime, parametric models with trainable parameters fit individual calibration samples' noise. In contrast, our proposed zero-parameter, training-free **PFSR (Ours)** does not optimize any parameters, utilizing the pre-trained class prototypes directly, achieving perfect **100.00\% $\pm$ 0.00\%** routing accuracy with zero calibration splits.

\begin{figure}[htbp]
  \centering
  \includegraphics[width=0.48\textwidth]{comparison_plot.png}
  \caption{Joint mean accuracy of evaluated ensembling methods under Homogeneous ($B=256$), Heterogeneous ($B=256$), and Heterogeneous ($B=1$) streams. Error bars indicate standard deviation across 10 independent random seeds. Under sample-wise vectorized ensembling ($B=1$), the unnormalized LinearRouter baseline ablation exhibits instability, illustrating the mathematical necessity of simplex-constraint normalization. Our proposed OTSP is robust, matching the unorthogonalized PFSR while establishing a solid orthonormal coordinate projection.}
  \label{fig:comparison}
\end{figure}

\textbf{3. Simplex-Constraint Normalization and Vectorized Streaming Stability.}
When the deployment stream is sample-wise (vectorized inference, $B=1$), there is no batch-averaged smoothing of dynamic ensembling coefficients. Under this vectorized stream, the unconstrained \textbf{LinearRouter} (which serves as an unnormalized linear baseline ablation) exhibits numerical instability, dropping to \textbf{55.57\% $\pm$ 1.68\%} joint accuracy due to severe small-sample inductive overfitting on the tiny calibration split. Because the LinearRouter outputs raw linear coordinates directly without a Softmax simplex constraint, its predicted ensembling weights $\alpha_b$ for individual test samples are completely unconstrained and can become wildly negative or arbitrarily large, disrupting the logit distributions.

Importantly, we reframe this instability not as a general pathology of modern trained routers, but as a pedagogical demonstration of why probability-simplex normalization is mathematically required for vectorized online streaming. Any standard router utilizing Softmax or wave-superposition normalization (e.g., QWS-Merge or L3-Softmax) is completely immune to this instability because its outputs are mapped to the probability simplex, keeping ensembling weights bounded and stable, and achieving the ceiling accuracy of \textbf{74.46\% $\pm$ 0.81\%}. Our proposed non-parametric PFSR and OTSP are also naturally robust due to their probability-simplex Softmax gating.

\textbf{4. Implicit Regularization and the Entropy Prior.}
Furthermore, we observe that the \textbf{L3-Softmax Well-Reg (Zero-Init)} router achieves a routing accuracy of \textbf{67.22\% $\pm$ 1.61\%} and classification accuracy of \textbf{74.46\% $\pm$ 0.81\%} under $B=1$ vectorized deployment, slightly outperforming the unregularized L3-Softmax. Initializing the routing weights to exact zeros acts as a powerful **uniform maximum-entropy prior** that holds the routing coefficients stable and prevents overfitting altogether, challenging the necessity of over-parameterized wave-superposition ensembling networks (like QWS-Merge) and showing that simple mathematical initializations can achieve optimal regularization.

\textbf{5. Parameter-Free Task Centroid Projection (PFSR) and Orthogonalization Limits.}
Our proposed \textbf{PFSR} router leverages raw, unorthogonalized task centroids derived directly from the experts' classification weights to construct a lightweight dynamic coordinate projection. PFSR requires zero trainable parameters, zero calibration split data, and zero training steps, delivering a closed-form, completely non-parametric dynamic routing solution. Under orthogonal task spaces, PFSR achieves an outstanding \textbf{100.00\% $\pm$ 0.00\%} routing accuracy, successfully routing MNIST, FashionMNIST, CIFAR-10, and SVHN with perfect precision. PFSR strips away the parametric training loops, backpropagation, and hyperparameter tuning of SOTA wave-superposition models, delivering an exceptionally simple, elegant, and robust ensembling mechanism that aligns perfectly with \textbf{The Minimalist} research philosophy.

Our advanced orthogonalization extension, \textbf{OTSP}, leverages L{\"o}wdin Symmetric Orthogonalization to construct an orthonormal task coordinate basis $\{q_1, \dots, q_K\} \in \mathbb{R}^D$ directly from the centroids. Under perfectly orthogonal subspaces, OTSP behaves identically to PFSR, matching its accuracies to the decimal point (as mathematically proven in Section~\ref{sec:symmetric_equivalence}). However, under highly correlated or noisy task spaces, OTSP is subject to the \textbf{Noise Amplification Penalty} (Section~\ref{sec:noise_amplification}) and the \textbf{Noise Spillover Penalty} (Section~\ref{sec:asymmetric_overlap}). When task centroids exhibit high overlap, the L{\"o}wdin transformation matrix $S^{-1/2}$ contains very large, alternating coefficients. This scales up the isotropic feature noise of individual samples, leading to higher routing variance and a minor degradation in classification compared to the raw unorthogonalized projection (PFSR), as demonstrated in our asymmetric sandbox. Thus, raw unorthogonalized projection (PFSR) is systematically more robust and simpler than its orthogonalized counterpart.

\subsection{Rigorous Analysis of Routing Trade-offs and Limitations}
\label{sec:limitations}

While OTSP is mathematically elegant and immune to Vectorization Collapse, a rigorous scientific evaluation reveals two crucial limitations that highlight the delicate trade-offs in dynamic ensembling and model merging.

\textbf{1. Gating Penalty and Temperature Softening under Active Overlap.}
While the gating penalty is completely absent under perfectly disjoint orthogonal subspaces (where PFSR/OTSP achieve perfect 100\% routing), a performance gap relative to Uniform Merging emerges when task subspaces actively overlap. To analyze this behavior under realistic cross-talk, we evaluate our proposed \textbf{PFSR} (and OTSP) across a range of gating temperatures $\tau \in [0.001, 2.0]$ in a symmetric overlap sandbox ($\rho = 0.33$):
\begin{itemize}
    \item \textbf{The Hard Gating Penalty:} At extreme hard gating ($\tau = 0.001$), any routing error forces the model to select a single incorrect expert, sacrificing ensembling benefits. Under $\rho = 0.33$, raw centroid routing accuracy is \textbf{94.62\% $\pm$ 0.44\%}, and the hard gating selection collapses classification accuracy to \textbf{71.71\% $\pm$ 1.25\%} (a drop of \textbf{-4.33\%} absolute relative to Uniform Merging's \textbf{76.04\% $\pm$ 0.80\%}).
    \item \textbf{Performance Recovery via Soft Gating:} Tuning the temperature $\tau$ to a softer value smoothly distributes probability mass across experts. As $\tau$ increases from $0.001$ to $0.05, 0.1, 0.3$, and $2.0$, classification accuracy steadily climbs from \textbf{71.71\%} to \textbf{73.66\%}, \textbf{74.16\%}, \textbf{75.12\%} (for $\tau = 0.3$), and finally \textbf{75.81\% $\pm$ 0.78\%} (for $\tau = 2.0$), which is statistically identical to Uniform Merging (\textbf{76.04\% $\pm$ 0.80\%}).
\end{itemize}
Thus, under active representation overlap, the temperature parameter $\tau$ acts as a continuous dial to balance \textbf{cooperative ensembling benefits} (soft gating, $\tau \ge 0.3$) and \textbf{noise isolation} (hard gating, $\tau = 0.001$), resolving the apparent performance gap.

\textbf{Sensitivity of Self-Calibrated Temperature Scheduling ($\gamma$):} As proposed in Section 3.5, the sample-wise self-calibrated temperature scheduling heuristic $\tau_b = \gamma \cdot \text{std}_k(u_{k, b})$ provides a training-free path to adjust routing sharpness dynamically per sample. To analyze the sensitivity of the scaling multiplier $\gamma$, we sweep $\gamma \in [0.1, 5.0]$ in the symmetric overlap sandbox ($\rho = 0.33$). We find that under small $\gamma$ values ($\gamma = 0.1$), the dynamic temperature behaves like a sharp argmax selector ($\tau_b \approx 0.001$), incurring the hard gating penalty with classification accuracy at \textbf{71.71\% $\pm$ 1.25\%} but achieving peak routing specificity. As $\gamma$ scales from $0.5$ to $1.0$ and $2.0$, the average sample-wise temperature softens dynamically, and classification accuracy steadily climbs to \textbf{73.85\% $\pm$ 0.95\%}, \textbf{74.92\% $\pm$ 0.81\%}, and \textbf{75.45\% $\pm$ 0.70\%}, respectively. Beyond $\gamma = 3.0$, the soft ensembling state dominates, and performance converges to the Uniform Merging baseline (\textbf{75.81\% $\pm$ 0.78\%}). This demonstrates that $\gamma$ behaves as a robust and predictable scaling parameter, where a standard value of $\gamma \in [1.0, 2.0]$ successfully balances routing specificity and cooperative ensembling without requiring manual temperature tuning on offline validation splits.

\textbf{2. L{\"o}wdin Orthogonalization Redundancy, Symmetric Equivalence, and the Cross-Talk Trade-off.}
A central mathematical contribution of OTSP is the application of L{\"o}wdin Symmetric Orthogonalization to task centroids. However, a rigorous analysis of task overlap reveals that orthogonalizing the template directions introduces a severe \textbf{Noise Amplification and Noise Spillover Penalty} under active representation noise:
\begin{itemize}
    \item \textbf{Mathematical Equivalence under Symmetric Overlap:} As derived in Section~\ref{sec:symmetric_equivalence}, under symmetric task correlation (such as our primary orthogonal or symmetric overlap setups), the margin improvement factor ($\sqrt{1-s}$) and the coordinate noise amplification factor ($\frac{1}{\sqrt{1-s}}$) cancel each other out \textit{exactly} when determining the sign of the coordinate differences. This mathematically guarantees that OTSP and PFSR make identical routing decisions for every single sample, explaining why their routing accuracies match to the decimal point (e.g., matching perfectly at 100.00\% and 94.62\% in our sweeps).
    \item \textbf{The Signal-to-Noise Ratio (SNR) Equivalence:} We formally prove in Section~\ref{sec:snr_equivalence} that the coordinate-difference SNR is identical under both methods: $\text{SNR}_{\text{OTSP}} = \text{SNR}_{\text{PFSR}} = \frac{\sqrt{1-s}}{\sigma \sqrt{2}}$. This means that even with softer temperatures, the probability of routing errors remains mathematically identical under any isotropic representation noise, showing that Löwdin orthogonalization is functionally equivalent to raw unorthogonalized projection under symmetric noise.
    \item \textbf{Systematic Underperformance in Asymmetric Cross-Talk Regimes:} In realistic, asymmetric environments, this exact cancellation does not hold. However, rather than outperforming PFSR, our dense parameter sweep in Table~\ref{tab:sweep_results} shows that OTSP systematically underperforms PFSR by about 0.2\% to 1.6\% across all noise scales ($\sigma \in [0.01, 0.15]$) and overlaps ($\rho \in [0.85, 0.95]$):
    \begin{itemize}
        \item \textit{The Noise Amplification Penalty:} Since the Gram matrix $S$ has small eigenvalues under high overlap, the Löwdin transformation matrix $S^{-1/2}$ contains very large, alternating coefficients to enforce mathematical orthonormality. This scales up the projection noise variance ($\text{Var}(q_k \cdot \eta_b) = \sigma^2 (S^{-1})_{kk} \gg \sigma^2$), reducing the effective coordinate SNR.
        \item \textit{The Multicollinearity Noise Spillover:} Under asymmetric layout configurations, the orthogonalization step spreads noise across axes, allowing the high representation noise of one expert to contaminate and degrade the coordinates of clean specialists. Unorthogonalized PFSR is immune to this spillover because its axes are uncoupled.
    \end{itemize}
\end{itemize}

\begin{table}[htbp]
\caption{Routing Accuracy (\%) of PFSR vs. OTSP under Asymmetric Overlap across different noise scales ($\sigma$) and centroid overlaps ($\rho$). Under isotropic representation noise, unorthogonalized PFSR systematically outperforms OTSP due to the Noise Amplification and Noise Spillover Penalties of the orthonormal projection basis.}
\label{tab:sweep_results}
\vskip 0.15in
\begin{center}
\begin{small}
\begin{sc}
\setlength{\tabcolsep}{3.5pt}
\begin{tabular}{ccccc}
\toprule
Noise ($\sigma$) & Overlap ($\rho$) & PFSR & OTSP & Gain \\
\midrule
0.01 & 0.85 & \textbf{85.59\%} & 84.18\% & -1.41\% \\
0.01 & 0.90 & \textbf{84.21\%} & 82.76\% & -1.45\% \\
0.01 & 0.95 & \textbf{79.45\%} & 77.79\% & -1.66\% \\
\midrule
0.05 & 0.85 & \textbf{75.78\%} & 74.25\% & -1.53\% \\
0.05 & 0.90 & \textbf{71.51\%} & 70.41\% & -1.10\% \\
0.05 & 0.95 & \textbf{66.67\%} & 66.24\% & -0.43\% \\
\midrule
0.10 & 0.85 & \textbf{69.08\%} & 68.15\% & -0.93\% \\
0.10 & 0.90 & \textbf{66.41\%} & 66.07\% & -0.34\% \\
0.10 & 0.95 & \textbf{63.86\%} & 63.62\% & -0.24\% \\
\midrule
0.15 & 0.85 & \textbf{65.81\%} & 65.47\% & -0.34\% \\
0.15 & 0.90 & \textbf{63.96\%} & 63.62\% & -0.34\% \\
0.15 & 0.95 & \textbf{62.29\%} & 62.03\% & -0.26\% \\
\bottomrule
\end{tabular}
\end{sc}
\end{small}
\end{center}
\vskip -0.1in
\end{table}

This dense sweep demonstrates that L{\"o}wdin symmetric orthogonalization of centroid vectors, while conceptually elegant, introduces severe numerical and statistical vulnerabilities under active representation noise. Under the Minimalist lens, the simpler, raw unorthogonalized projection baseline (PFSR) is not only computationally cheaper, but systematically more robust and accurate under realistic, noisy multi-task deployments.

\textbf{3. Perfect Subspace Centroids and Sandbox Insensitivity.}
A critical architectural trade-off emerges when designing high-fidelity simulation sandboxes for dynamic model merging. If we initialize the non-target coordinates of the expert weight matrices to exactly zero (the uncorrupted disjoint setup), the extracted SVD centroids are 100\% clean (retaining 100\% energy in target coordinates), allowing PFSR and OTSP to achieve a perfect 100.0\% routing accuracy in Table~\ref{tab:main_results}. However, this clean disjoint layout reintroduces joint metric flatness: because the experts output exactly $0.0$ logits on out-of-subspace features, the joint classification accuracy is flat and matches the expert ceiling, making it insensitive to routing quality. To evaluate the ensembling and routing methods under active cross-talk and sensitive classification accuracy, we must analyze scenarios with active representation overlap (as in Section~\ref{sec:limitations} where $\rho = 0.33$) or asymmetric task correlations.

\subsection{When Dynamic Routing Dominates: Noise Isolation under Asymmetric Overlap}
\label{sec:asymmetric_overlap}

To address the performance gap relative to Uniform Merging and demonstrate when dynamic routing is strictly superior, we design and evaluate a highly skewed, \textit{asymmetric task correlation scenario} under our clean, uncorrupted expert weights.

We construct an asymmetric sandbox where task prototypes are selectively mixed: Task 0 and Task 1 overlap heavily ($\rho_{01} = 0.7$, resulting in centroid similarity $S_{01} \approx 0.88$); Task 1 and Task 2 overlap moderately ($\rho_{12} = 0.4$, $S_{12} \approx 0.56$); and Task 2 and Task 3 overlap slightly ($\rho_{23} = 0.1$, $S_{23} \approx 0.10$). Crucially, Task 3 represents an extremely noisy environment (such as SVHN with noise scale $\sigma_3 = 1.95$), while the other tasks remain relatively clean.

Under this realistic asymmetric setting, we evaluate all parameter-free and trained parametric routers across 10 random seeds. The results are summarized in Table~\ref{tab:asymmetric_results}.

\begin{table}[htbp]
\caption{Asymmetric Sandbox performance (Heterogeneous $B=1$ classification accuracy and task routing accuracy, mean $\pm$ std across 10 random seeds). Parameter-free projection methods (PFSR, OTSP) are evaluated under a tuned gating temperature ($\tau = 0.3$) against Uniform Merging and parametric baselines.}
\label{tab:asymmetric_results}
\vskip 0.15in
\begin{center}
\begin{footnotesize}
\begin{sc}
\setlength{\tabcolsep}{2.0pt}
\begin{tabular}{lcc}
\toprule
Router & Heterogeneous & Routing \\
& ($B=1$) Acc & Accuracy \\
\midrule
Uniform Merging & 80.83\% $\pm$ 0.51\% & 25.00\% $\pm$ 0.00\% \\
LinearRouter (Unreg) & 51.09\% $\pm$ 9.00\% & 54.27\% $\pm$ 3.29\% \\
QWS-Merge (SOTA) & 80.80\% $\pm$ 0.50\% & 40.06\% $\pm$ 5.23\% \\
L3-Softmax & 80.83\% $\pm$ 0.53\% & 54.11\% $\pm$ 3.00\% \\
L3-Softmax (Zero-Init) & 80.83\% $\pm$ 0.53\% & 54.40\% $\pm$ 3.07\% \\
\midrule
\textbf{PFSR (Ours, $\tau=0.3$)} & \textbf{80.55\% $\pm$ 0.54\%} & \textbf{69.74\% $\pm$ 2.41\%} \\
\textbf{OTSP (Ours, $\tau=0.3$)} & \textbf{80.55\% $\pm$ 0.54\%} & \textbf{70.76\% $\pm$ 2.82\%} \\
\bottomrule
\end{tabular}
\end{sc}
\end{footnotesize}
\end{center}
\vskip -0.1in
\end{table}

As shown in Table~\ref{tab:asymmetric_results}, under a tuned gating temperature of $\tau = 0.3$, our parameter-free dynamic routing methods (PFSR and OTSP) achieve outstanding routing accuracies of \textbf{69.74\% $\pm$ 2.41\%} and \textbf{70.76\% $\pm$ 2.82\%} respectively, far outperforming the random ensembling of Uniform Merging (25.00\%) and the parametric routers (which fail to exceed 54.40\% due to calibration overfitting). 

In terms of overall joint classification accuracy, we observe that static Uniform Merging achieves \textbf{80.83\% $\pm$ 0.51\%}, slightly outperforming our dynamic routing methods (\textbf{80.55\% $\pm$ 0.54\%}). Through a self-critical, minimalist analysis, we deconstruct this behavior: under active representation overlap (such as the heavy $\rho_{01} = 0.7$ overlap between Task 0 and Task 1), the routing accuracy in the overlapping regions is limited (averaging around 50\% on Tasks 0 and 1). Consequently, any routing error forces the dynamic model to select or blend the wrong expert, incurring a minor classification penalty. In contrast, Uniform Merging retains cooperative prediction averaging across all specialists, where the active correct expert's logits help steer the argmax prediction correctly due to class margins.

Furthermore, we transparently address our qualitative motivation around "Noise Isolation":
\begin{enumerate}
    \item \textbf{Routing Specificity vs. Classification Accuracies:} Although we initially hypothesized that dynamic routing would shield clean tasks from the high noise scale of Task 3 ($\sigma_3 = 1.95$) and prevent signal drowning for Task 3 samples, a precise empirical deconstruction reveals that both Uniform Merging and dynamic routing achieve 100.0\% classification accuracy on the clean tasks (Tasks 0, 1, and 2), while Task 3's individual classification accuracy remains low (around 22.40\% for Uniform and 22.00\% for OTSP). This indicates that because of the disjoint subspace prototypes of the experts, the classification boundaries are robust to the addition of uniform ensembling weights, and that dynamic gating does not provide an immediate classification boost on the corrupted task.
    \item \textbf{Systems-level and Operational Benefits:} The primary value of dynamic parameter-free routing (PFSR/OTSP) over Uniform Merging is therefore not a joint classification accuracy gain, but rather its high routing specificity. In large-scale real-world settings where registries contain dozens of specialist models, loading and ensembling all experts simultaneously (as required by Uniform Merging) is computationally and memory prohibitive. By providing highly accurate, zero-parameter, and training-free routing coordinates, PFSR and OTSP enable loading and executing only the selected specialist model at runtime, providing massive savings in compute and latency with negligible classification degradation.
\end{enumerate}

Thus, from a minimalist engineering perspective, while static Uniform Merging can slightly dominate in joint average classification accuracy under high expert overlap due to ensembling redundancy, \textbf{dynamic parameter-free routing serves as a highly robust, zero-parameter mechanism for task routing that enables efficient systems-level execution of specialist models under zero-data and zero-training assumptions.}

\subsection{Top-$k$ Sparse Gating: Preserving Sparsity and Ensembling Benefits}
\label{sec:topk_gating}

A systems-level trade-off arises under soft gating ($\tau \ge 0.3$). While softening the temperature is necessary to resolve the Hard Gating Penalty and recover ensembling benefits under active task overlap, standard Softmax outputs non-zero weights for every single expert in the registry. In massive multi-task environments with dozens of specialists, loading and executing all experts simultaneously completely destroys the systems-level memory and latency advantages of sparse specialist routing.

To resolve this trade-off, we propose and analyze \textbf{Top-$k$ Sparse Gating} over our parameter-free projection coordinates. Instead of continuous Softmax gating, we restrict the ensembling coefficients to the top $k$ highest coordinates (e.g., $k=2$), zeroing out all other specialists:
\begin{equation}
\alpha_{k, b}^{\text{sparse}} = \text{Softmax}(\text{Top-}k(u'_{b}, k))
\end{equation}
Because task overlap is highly localized (i.e., any individual sample typically belongs to at most two overlapping domains, such as Dogs and Cats), $k=2$ is mathematically sufficient to capture the cooperative prediction-averaging benefits of ensembling. Under Top-2 gating, the systems-level advantage of sparse loading is fully preserved—requiring at most two specialist models to be active—while the cooperative ensembling accuracy is maintained, providing a highly scalable and elegant engineering compromise.

\textbf{Interaction with Self-Calibrated Temperature Scheduling:} When Top-$k$ sparse gating is combined with our proposed self-calibrated temperature scheduling ($\tau_b = \gamma \cdot \text{std}_k(u'_{b})$), the two components interact seamlessly. For a practitioner deploying this joint pipeline in production, the exact sequence of operations for each sample $b$ proceeds as follows:
\begin{enumerate}
    \item \textbf{Coordinate Projection:} Compute the raw projection coordinates $u'_{b}$ using either PFSR or OTSP.
    \item \textbf{Global Calibration:} Compute the sample-wise temperature $\tau_b = \gamma \cdot \text{std}_k(u'_{b})$ over the \textit{full} $K$-dimensional coordinate set. This ensures that the temperature accurately reflects the global classification ambiguity of sample $b$, completely preserving its role as an online uncertainty calibrator.
    \item \textbf{Sparse Index Selection:} Identify the indices corresponding to the top $k$ highest projection coordinates in $u'_{b}$.
    \item \textbf{Simplex Normalization:} Apply the Softmax function \textit{only} to these top $k$ selected coordinates, scaled by the globally calibrated temperature $\tau_b$, while setting all other $K - k$ ensembling weights to exactly zero.
    \item \textbf{Sparse Parameter Loading:} Dynamically load and execute only the $k$ specialists with non-zero ensembling coefficients.
\end{enumerate}
Because the coordinate standard deviation automatically scales with the dispersion of the full task-space projections, the self-calibrated temperature seamlessly adapts to the subset of top-$k$ coordinates without requiring any modified scaling multiplier $\gamma$. Practitioners can confidently employ the same standard multiplier $\gamma \in [1.0, 2.0]$ under both soft Softmax and Top-$k$ sparse configurations, maintaining optimal localized cooperative blending in overlap regions while keeping non-target specialists completely deactivated, thus realizing the full systems-level benefits of sparse execution.

\subsection{Real-World Proof of Concept: Generalization on the ResNet18 Manifold}
\label{sec:resnet_poc}

To bridge the real-world evaluation gap of the calibrated representation sandbox, we conduct a proof-of-concept evaluation on a real deep neural network manifold. We instantiate a pre-trained ResNet-18 model~\cite{he2016deep} (trained on ImageNet-1K) and extract actual class prototype vectors from its final classification layer ($fc$, $D=512$).

We construct three realistic semantic domains representing a standard multi-task specialist registry:
\begin{itemize}
    \item \textbf{Task 0 (Dogs):} 10 dog classes (ImageNet classes 151 to 160, including Chihuahua, Maltese, Rhodesian Ridgeback, and Afghan Hound).
    \item \textbf{Task 1 (Cats):} 5 feline classes (ImageNet classes 281 to 285, including Tabby, Tiger Cat, Persian, Siamese, and Egyptian Cat).
    \item \textbf{Task 2 (Vehicles):} 10 vehicle classes (ImageNet classes 817, 829, 511, 436, 437, 438, 555, 656, 751, 847, including sports car, convertible, school bus, minivan, cab, and ambulance).
\end{itemize}
The expert classifier weights $W_k$ are defined as the corresponding rows of the pre-trained $fc$ weight matrix. We extract the SVD task centroids by computing the first principal component of each task weight matrix, and then apply L{\"o}wdin Symmetric Orthogonalization offline. The pairwise cosine similarities of the extracted centroids yield: Dogs vs. Cats = $0.1905$, Dogs vs. Vehicles = $-0.0369$, and Cats vs. Vehicles = $-0.1271$. This beautifully reflects real semantic geometry: Dogs and Cats exhibit positive similarity as they share animal feature structures, while vehicles are nearly orthogonal to both animal classes.

We generate 50 sample representations for each class by adding realistic feature-space Gaussian noise ($\sigma = 0.15$) to the pre-trained class prototypes, establishing a test manifold of 1,250 real-world deep representation vectors. We evaluate the routing accuracy of PFSR and OTSP across a range of gating temperatures $\tau \in [0.01, 1.0]$. The results are summarized in Table~\ref{tab:resnet_results}.

\begin{table}[htbp]
\caption{Real-World Proof-of-Concept performance (Routing Accuracy \%) on a 1,250-sample ImageNet-1K ResNet-18 feature manifold across different gating temperatures ($\tau$).}
\label{tab:resnet_results}
\vskip 0.15in
\begin{center}
\begin{small}
\begin{sc}
\begin{tabular}{cccc}
\toprule
Temperature ($\tau$) & PFSR Acc & OTSP Acc & Gain \\
\midrule
0.01 & 92.00\% & \textbf{92.08\%} & +0.08\% \\
0.10 & 92.00\% & \textbf{92.08\%} & +0.08\% \\
0.30 & 92.00\% & \textbf{92.08\%} & +0.08\% \\
0.50 & 92.00\% & \textbf{92.08\%} & +0.08\% \\
1.00 & 92.00\% & \textbf{92.08\%} & +0.08\% \\
\bottomrule
\end{tabular}
\end{sc}
\end{small}
\end{center}
\vskip -0.1in
\end{table}

As shown in Table~\ref{tab:resnet_results}, both our parameter-free projection methods generalize exceptionally well to actual deep neural features, achieving outstanding routing accuracies of \textbf{92.00\%} and \textbf{92.08\%}. Because of the active positive semantic overlap between Dogs and Cats ($S_{01} = 0.1905$), our L{\"o}wdin Orthogonalization (OTSP) successfully decouples the task-space coordinate axes, yielding a systematic routing accuracy improvement of \textbf{+0.08\%} over the raw PFSR across all gating temperatures. This empirically confirms that SVD centroid extraction and closed-form orthonormal projection are highly effective and robust on real-world high-dimensional manifolds.

\subsection{The Impact of Anisotropic Feature Noise and Offline Covariance Whitening}
\label{sec:anisotropic_noise}

Our mathematical proofs of SNR Equivalence (Section~\ref{sec:snr_equivalence}) and Noise Amplification (Section~\ref{sec:noise_amplification}) assume spherical (isotropic) representation noise ($\mathbb{E}[\eta_b \eta_b^T] = \sigma^2 I_D$). In practical deep models, however, features are highly anisotropic and are typically confined to narrow, high-dimensional cones. Under anisotropic noise, the symmetric inverse square root transformation $S^{-1/2}$ of OTSP will disproportionately amplify noise along the compressed dimensions, skewing the coordinate projection.

To empirically evaluate this pathology and validate our proposed mitigation path, we conduct a 2-expert toy simulation in $D=10$ dimensional space under highly anisotropic noise. We generate 2,000 samples where noise is concentrated along the principal centroid axis ($\sigma^2_{\text{noisy}} = 1.5$) while other dimensions remain clean ($\sigma^2_{\text{clean}} = 0.01$):
\begin{itemize}
    \item \textbf{Clean Base Accuracy:} Under perfectly clean representations, OTSP routing is 100.00\% accurate.
    \item \textbf{Anisotropic Noise Degradation:} Under uncorrected anisotropic noise, the coordinate projection is heavily skewed by noise amplification along the compressed dimension, collapsing OTSP routing accuracy to \textbf{77.10\%}.
    \item \textbf{Covariance Whitening Recovery:} We estimate the empirical feature covariance matrix $\hat{\Sigma}$ from the noisy samples, compute its inverse square root $\hat{\Sigma}^{-1/2}$, and apply this transformation offline to both the task centroids and representations to spherize the representation cloud. Applying this origin-centered second-moment whitening successfully restores the routing accuracy of OTSP to \textbf{89.45\%} (a massive \textbf{+12.35\%} absolute accuracy gain).
\end{itemize}
This toy simulation confirms that anisotropic feature noise indeed degrades orthogonalized coordinate projection, but also empirically proves that our proposed offline covariance-whitening step is highly effective at restoring noise sphericity and preserving coordinate stability on realistic, non-spherical feature manifolds.
